\section{Results}
\label{sec:results}

We present evidence for the ordering effect (RQ1) and its mechanism (RQ2).

\subsection{Main Results (RQ1)}

\begin{table}[h]
\centering
\caption{Main results on HumanEval + MBPP}
\label{tab:main_results}
\begin{tabular}{@{}lcc@{}}
\toprule
Condition & pass@1 & 95\% CI \\
\midrule
Static$\rightarrow$Execution (A) & \textbf{55.57\%} & --- \\
Execution$\rightarrow$Static (B) & 43.07\% & --- \\
\midrule
\textbf{Relative Improvement} & \textbf{29.02\%} & [15.74\%, 44.53\%] \\
McNemar $p$-value & $6.31\times10^{-6}$ & --- \\
\bottomrule
\end{tabular}
\end{table}

\textbf{Key Observations:}

\begin{enumerate}
    \item \textbf{Large effect size:} Static$\rightarrow$execution ordering achieves 29.02\% relative improvement, nearly doubling our 15\% threshold and substantially exceeding prior self-repair improvements.

    \item \textbf{Strong statistical significance:} $p = 6.31\times10^{-6}$, with 95\% CI excluding zero by a wide margin.

    \item \textbf{Matched content control:} The 12.5 percentage point absolute improvement stems purely from presentation order.
\end{enumerate}

\subsection{Mechanism Analysis (RQ2)}

\begin{table}[h]
\centering
\caption{Early iteration gains (H-M1)}
\label{tab:early_gains}
\begin{tabular}{@{}lcc@{}}
\toprule
Condition & $\Delta$Pass$_{1\rightarrow2}$ & $p$-value \\
\midrule
Static-First (A) & 12.50\% & --- \\
Exec-First (B) & 12.05\% & --- \\
\midrule
Difference & +0.45pp & 0.859 \\
\bottomrule
\end{tabular}
\end{table}

\textbf{Finding:} Early gains are nearly identical. H-M1 is \textbf{not supported}.

\begin{table}[h]
\centering
\caption{Regression rates (H-M2)}
\label{tab:regression}
\begin{tabular}{@{}lcc@{}}
\toprule
Condition & Regression Rate$_{1\rightarrow2}$ & $p$-value \\
\midrule
Static-First (A) & \textbf{21.53\%} & --- \\
Exec-First (B) & 34.69\% & --- \\
\midrule
Difference & 13.15pp (38\%) & 0.0198 \\
\bottomrule
\end{tabular}
\end{table}

\textbf{Finding:} Static-first ordering reduces regressions by 38\%. H-M2 is \textbf{supported}.

\subsection{Summary}

The effect operates through trajectory stabilization, not acceleration. Static-first ordering helps the LLM repair more stably, not faster.
